% Copia mecánica íntegra de apendices.tex:1--761.
% Residencia material del ensamblaje sucesor de 102 capítulos.
\part{Appendices démonstratifs et tableaux de classification}

\chapter{Calcul affine ordonné des feuillets APP}

\section{Produit semi-direct affine et composition des voies}

Soit \(U_{\APP}\) le module radial antérieur à son évaluation archimédienne. L'enveloppe correcte est le monoïde
\[
 \AffEnd(U_{\APP})
 =\operatorname{End}(U_{\APP})\ltimes U_{\APP},
\]
car les feuillets de seuil ou de pliage peuvent contenir des endomorphismes non inversibles. L'action de \((\Lambda,C)\) sur \(u\in U_{\APP}\) est
\[
 (\Lambda,C)\cdot u=\Lambda u+C.
\]
La composition se calcule sans commuter les facteurs :
\begin{align}
 (\Lambda_2,C_2)(\Lambda_1,C_1)\cdot u
 &=\Lambda_2(\Lambda_1u+C_1)+C_2 \notag\\
 &=(\Lambda_2\Lambda_1)u+(C_2+\Lambda_2C_1).
 \label{espiral:eq:ap-comp-afin}
\end{align}
Par conséquent,
\[
 (\Lambda_2,C_2)(\Lambda_1,C_1)
 =(\Lambda_2\Lambda_1,C_2+\Lambda_2C_1).
\]
La seconde coordonnée n'est pas une somme ordinaire de déplacements : c'est un cocycle tordu par l'action linéaire accumulée.

\begin{proposition}[Formule des préfixes]
Pour une voie de \(N\) facteurs \(h_j=(\lambda_j,C_j)\), ordonnés de \(j=0\) à \(j=N-1\), l'holonomie préfixe est
\[
 H_N=h_{N-1}\cdots h_1h_0=(\Lambda_N,\mathcal C_N),
\]
où
\[
 \Lambda_N=\prod_{j=0}^{N-1}\lambda_j
\]
et
\[
 \mathcal C_N
 =C_{N-1}+\lambda_{N-1}C_{N-2}
 +\lambda_{N-1}\lambda_{N-2}C_{N-3}
 +\cdots+
 \lambda_{N-1}\cdots\lambda_1C_0.
\]
\end{proposition}
\begin{proof}
Pour \(N=1\), la formule est immédiate. Si elle est valable pour \(N\), alors
\[
 h_NH_N
 =(\lambda_N\Lambda_N,C_N+\lambda_N\mathcal C_N),
\]
ce qui produit exactement l'expression de longueur \(N+1\). La récurrence préserve l'ordre de la voie et n'identifie pas les parcours ayant les mêmes extrémités.
\end{proof}

\section{Différence somme--produit comme courbure d'ordre}

Soient \(A_a=(I,a)\) et \(M_\lambda=(\lambda,0)\). Dans le secteur inversible,
\[
 A_aM_\lambda A_a^{-1}M_\lambda^{-1}
 =A_{a(1-\lambda)}.
\]
La démonstration directe sur \(u\) est
\[
 u\xmapsto{M_\lambda^{-1}}\lambda^{-1}u
 \xmapsto{A_a^{-1}}\lambda^{-1}u-a
 \xmapsto{M_\lambda}u-\lambda a
 \xmapsto{A_a}u+a(1-\lambda).
\]
Ainsi, la différence entre accumuler et plier est publiée sous forme de translation. Lorsque \(\lambda=1\), le commutateur s'annule ; lorsque \(\lambda\neq1\), la voie conserve une holonomie qui ne peut être reconstruite à partir du seul point final.

\begin{example}[Deux voies avec la même donnée visible]
Fixons \(u_0=0\), \(a=1\), \(\lambda=2\). Les voies
\[
 A_1M_2:u_0\longmapsto1,
 \qquad
 M_2A_{1/2}:u_0\longmapsto1
\]
publient la même valeur finale. Cependant, leurs registres contiennent des facteurs, déplacements et préfixes distincts. L'évaluation ponctuelle identifie les voies ; l'holonomie affine enrichie les sépare.
\end{example}

\section{Tableau exhaustif de la carte locale à 2 orientations}

Le tableau suivant ne postule pas que tout degré radial global provient d'une seule paire. Il enregistre la carte élémentaire obtenue en composant deux orientations TRIT. Le degré d'une voie générale est le cocycle entier de \cref{espiral:eq:cociclos-radiales} ; cette carte fournit ses neuf générateurs locaux.

\begin{longtable}{rrrrrrp{.27\textwidth}}
\caption{Les \(9\) paires orientées de la carte TRIT locale.}
\label{espiral:tab:ap-nueve-pares}\\
\toprule
\(\tau_+\)&\(\tau_-\)&\(p\)&\(a\)&\(r_3\)&\(c_3\)&Lectura local\\
\midrule
\endfirsthead
\toprule
\(\tau_+\)&\(\tau_-\)&\(p\)&\(a\)&\(r_3\)&\(c_3\)&Lectura local\\
\midrule
\endhead
\( 1\)&\( 1\)&\( 2\)&\( 0\)&\(-1\)&\( 1\)&ouverture radiale double avec retenue positive\\
\( 1\)&\( 0\)&\( 1\)&\( 1\)&\( 1\)&\( 0\)&ouverture et seuil\\
\( 1\)&\(-1\)&\( 0\)&\( 2\)&\( 0\)&\( 0\)&équilibre à orientation positive\\
\( 0\)&\( 1\)&\( 1\)&\(-1\)&\( 1\)&\( 0\)&seuil et ouverture\\
\( 0\)&\( 0\)&\( 0\)&\( 0\)&\( 0\)&\( 0\)&umbral doble\\
\( 0\)&\(-1\)&\(-1\)&\( 1\)&\(-1\)&\( 0\)&seuil et contraction\\
\(-1\)&\( 1\)&\( 0\)&\(-2\)&\( 0\)&\( 0\)&équilibre à orientation négative\\
\(-1\)&\( 0\)&\(-1\)&\(-1\)&\(-1\)&\( 0\)&contraction et seuil\\
\(-1\)&\(-1\)&\(-2\)&\( 0\)&\( 1\)&\(-1\)&contraction double avec retenue négative\\
\bottomrule
\end{longtable}

Ici \(p=r_3+3c_3\), avec le résidu équilibré \(r_3\in\{-1,0,1\}\). La retenue distingue \(p=2\) de son résidu \(-1\) et \(p=-2\) de son résidu \(1\). La paire \((p,a)\) distingue également les trois généalogies locales de degré \(0\).

\chapter{Certificat nonadique et géométrie de la mémoire}

\section{Distinction typée entre catalogue décimal, mots ternaires et fibre ambiante}

\begin{longtable}{p{.18\textwidth}p{.18\textwidth}p{.52\textwidth}}
\caption{Strates du certificat de la connexion nonadique conjointe.}
\label{espiral:tab:certificado-nonadico}\\
\toprule
Objet&Cardinalité&Fonction mathématique\\
\midrule
\endfirsthead
\toprule
Objet&Cardinalité&Fonction mathématique\\
\midrule
\endhead
Catalogue des triplets d'origine&\(104\,976\)&Domaine fini antérieur à la sélection conjointe ; ce n'est pas un tableau décimal.\\
Émissions décimales \(U_6\)&\(468\)&Émissions compatibles de longueur \(6\) dans la publication décimale.\\
Mots visibles \(w_6\)&\(243\)&Mots ternaires publiés après le typage ; ils ne constituent pas l'espace ambiant complet.\\
Espace ambiant \(\F_3^6\)&\(729\)&Espace de tous les mots ternaires de longueur \(6\) ; il contient les \(243\) visibles.\\
Niveaux certifiés&\(396\)&Prolongement fini jusqu'aux bords déclarés.\\
Trits recorridos&\(2\,376\)&Longitud total \(396\cdot6\).\\
Vueltas completas&\(43\)&Cycles nonadiques complets dans le segment certifié.\\
Paires \((K,K+9)\)&\(387\)&Paires de même phase par canal ; le registre peut différer.\\
\bottomrule
\end{longtable}

La chaîne de prolongement se conserve dans l'ordre
\[
 w_6\longrightarrow w_{12}\longrightarrow w_{18}
 \longrightarrow w_{24}\longrightarrow w_{30}
 \longrightarrow R_{36}\longrightarrow G_9.
\]
Aucun de ces objets ne peut être remplacé par le quotient modal ultérieur. L'holonomie \(\Gamma_9\) agit sur la dynamique enrichie ; la publication de mémoire réduite ne conserve que l'information spécifiée par son codomaine.

\section{Retour de phase, avancée de mémoire et état complet}

Avec
\[
 \rho_9(N)=1+((N-1)\bmod9),
 \qquad
 q_9(N)=\left\lfloor\frac{N-1}{9}\right\rfloor,
\]
on a
\[
 \rho_9(N+9)=\rho_9(N),
 \qquad q_9(N+9)=q_9(N)+1.
\]
Par conséquent,
\[
 H_9(N+9)=H_9(N)+(0,0,1).
\]
L'égalité de la première coordonnée de phase n'implique pas l'égalité de l'état :
\[
 g(K+9)=g(K),
 \qquad B_{K+9}\neq B_K
\]
en général. La formule hélicoïdale est précisément la représentation minimale de cette distinction.

\section{Extension dodécaphasique non scindée}

La suite
\[
 0\longrightarrow C_{12}\longrightarrow C_{108}
 \longrightarrow C_9\longrightarrow0
\]
couple le calendrier dodécaphasique au retour nonadique. S'il existait une section homomorphe globale \(s:C_9\to C_{108}\), le générateur de \(C_9\) devrait se relever en un élément d'ordre divisant \(9\) qui se projette sur le générateur. Le relèvement temporel actif conserve cependant le déplacement dodécaphasique accumulé et n'identifie pas l'état après un tour de phase. La couture est enregistrée sous forme de retenue ; géométriquement, elle se manifeste par une translation axiale de l'hélice.

\section{5 régions de \texorpdfstring{\(\pi\)}{pi} et non-fidélité de l'évaluation}

La microfibre de \(\pi\) contient cinq régions orientées de multiplicité totale
\[
 432+4\cdot144=1008.
\]
Les régions publient la même valeur réelle, mais conservent des généalogies distinctes. Ce résultat fournit le précédent interne de la classification spirale : une courbe visible ne détermine pas à elle seule la région, le feuillet, la retenue ni la mémoire qui l'ont produite.

\chapter{Algèbre puissance--logarithmique et 5 publications canoniques}

\section[Coordonnée puissance--logarithmique et inverse]{La coordonnée \(L_p\) et son inverse}

Pour \(x>0\), on définit
\[
 L_p(x)=
 \begin{cases}
 (x^p-1)/p,&p\neq0,\\
 \log x,&p=0.
 \end{cases}
\]
Son inverse est
\[
 \exp_p(u)=
 \begin{cases}
 (1+pu)^{1/p},&p\neq0,\ 1+pu>0,\\
 \e^u,&p=0.
 \end{cases}
\]
La dérivée
\[
 \frac{\dd}{\dd x}L_p(x)=x^{p-1}>0
\]
prouve que \(L_p\) est strictement croissante sur son domaine et donc inversible.

\begin{proposition}[Loi somme--produit]
Pour \(x,y>0\),
\[
 L_p(xy)=L_p(x)+L_p(y)+pL_p(x)L_p(y).
\]
\end{proposition}
\begin{proof}
Si \(p\neq0\),
\begin{align*}
 L_p(x)+L_p(y)+pL_p(x)L_p(y)
 &=\frac{x^p-1}{p}+\frac{y^p-1}{p}
   +\frac{(x^p-1)(y^p-1)}p\\
 &=\frac{x^py^p-1}{p}=L_p(xy).
\end{align*}
Pour \(p=0\), l'égalité se réduit à \(\log(xy)=\log x+\log y\).
\end{proof}

\begin{corollary}[Groupe déformé sur le domaine positif]
L'opération
\[
 u\boxplus_pv=u+v+puv
\]
satisfait
\[
 1+p(u\boxplus_pv)=(1+pu)(1+pv).
\]
Son neutre est \(0\) et l'inverse de \(u\) est
\[
 \ominus_pu=-\frac{u}{1+pu},
\]
si \(1+pu\neq0\).
\end{corollary}

\section{Dérivation des 5 courbes canoniques}

Soit \(m\) la profondeur publiée et supposons
\[
 L_p(r/r_0)=\beta m,
 \qquad
 \theta=\theta_0+\omega m,
 \qquad\omega\neq0.
\]
Eliminando \(m=(\theta-\theta_0)/\omega\),
\[
 r(\theta)=r_0
 \left[1+p\frac{\beta}{\omega}
 (\theta-\theta_0)\right]^{1/p},
 \qquad p\neq0,
\]
et
\[
 r(\theta)=r_0
 \exp\!\left[\frac{\beta}{\omega}
 (\theta-\theta_0)\right]
\]
pour \(p=0\). Les changements d'origine et les changements d'échelle positifs produisent les formes conventionnelles :

\begin{longtable}{rllp{.31\textwidth}}
\caption{Publications canoniques du secteur translationnel.}
\label{espiral:tab:cinco-publicaciones-detalle}\\
\toprule
\(p\)&Équation normale&Nom conventionnel&Domaine et caractéristique\\
\midrule
\endfirsthead
\toprule
\(p\)&Équation normale&Nom conventionnel&Domaine et caractéristique\\
\midrule
\endhead
\(2\)&\(r^2=a^2\theta\)&Fermat&\(\theta\ge0\) ; accroissement linéaire de l'aire radiale quadratique.\\
\(1\)&\(r=a\theta\)&Archimède&\(\theta\ge0\) ; accroissement radial constant par unité angulaire.\\
\(0\)&\(r=r_0\e^{b\theta}\)&logarithmique&\(\theta\in\R\) ; dilatation constante par unité angulaire.\\
\(-1\)&\(r=a/\theta\)&hyperbolique&\(\theta>0\) ; produit \(r\theta\) constant.\\
\(-2\)&\(r^2=a^2/\theta\)&lituus&\(\theta>0\) ; produit \(r^2\theta\) constant.\\
\bottomrule
\end{longtable}

\section{Inversion des feuillets}

Para \(p\neq0\),
\[
 L_{-p}(x^{-1})
 =\frac{x^p-1}{-p}=-L_p(x).
\]
L'identité vaut aussi pour \(p=0\). C'est pourquoi l'inversion radiale échange \(p\) avec \(-p\) et change le signe de la coordonnée linéarisée. Cette relation regroupe Fermat avec le lituus et Archimède avec l'hyperbolique comme orientations conjuguées du lecteur, sans les identifier comme courbes égales.

\section{Secteur affine général}

Si la coordonnée radiale satisfait
\[
 u_{n+1}=\Lambda u_n+C,
\]
l'itération interne exacte est
\[
 u_n=\Lambda^n u_0+\sum_{k=0}^{n-1}\Lambda^kC.
\]
Si \(I-\Lambda\) est inversible, \(u_*=(I-\Lambda)^{-1}C\) et \(u_n-u_*=\Lambda^n(u_0-u_*)\). Ce n'est que dans une réalisation scalaire positive, \(\chi\Lambda=\lambda\chi\) avec \(\lambda>0\), que l'on définit \(u_*=C/(1-\lambda)\). Avec \(\theta_n=\theta_0+n\omega\), \(\omega\neq0\), l'interpolation publiée dans cette réalisation est
\[
 r(\theta)=r_0\exp_p\!\left[
 u_*+\lambda^{(\theta-\theta_0)/\omega}(u_0-u_*)
 \right].
\]
La taxonomie complète n'est donc pas une liste de cinq noms. Elle contient au moins le degré \(p\), l'orientation \(a\), la classe affine \([\Lambda,C]\), le cocycle de retenue, le bord, l'incrément angulaire et le registre de mémoire.

\chapter{Naturalité, équivariance et non-fidélité}

\section{Le lecteur de préfixes}

Pour une histoire tronquée
\[
 \widehat\gamma_N=\widehat e_N\cdots\widehat e_1,
\]
le lecteur enrichi enregistre
\[
 \mathcal P_{\rad,N}^{\enr}(\widehat\gamma_N)
 =
 \left(
 (H_0,\Led_0);
 \bigl(p(\widehat e_j),a(\widehat e_j),P_j,A_j,
 h_{e_j},H_j,\Led_j\bigr)_{j=1}^{N}
 \right).
\]
Le terme \(H_j=h_{e_j}\cdots h_{e_1}\) conserve les préfixes. Le registre \(\Led_j\) conserve, selon la voie, le feuillet, le résidu, le quotient, la retenue, le bord, l'orientation, le cylindre, la signature, la survie et la mémoire. Les coordonnées réelles s'obtiennent ensuite :
\[
 \mathcal P_{\arq}^{\mathrm{esp}}
 =\ev_{\R^2}\circ\mathcal P_{\rad}^{\enr}.
\]

\section{Preuve complète de naturalité sous troncature}

Soit \(\pi_{N,M}\) l'effacement du suffixe de longueur \(N-M\), avec \(M<N\). Soit \(\operatorname{pr}_{N,M}\) la projection qui conserve l'enregistrement initial et les \(M\) premiers enregistrements d'arête. Alors
\[
 \operatorname{pr}_{N,M}\,
 \mathcal P_{\rad,N}^{\enr}(\widehat\gamma_N)
 =
 \left(
 (H_0,\Led_0);
 \bigl(p(\widehat e_j),a(\widehat e_j),P_j,A_j,
 h_{e_j},H_j,\Led_j\bigr)_{j=1}^{M}
 \right).
\]
Par définition de la troncature,
\[
 \pi_{N,M}(\widehat\gamma_N)=\widehat e_M\cdots\widehat e_1,
\]
et en appliquant de nouveau le lecteur, on obtient la même séquence. Par conséquent,
\[
 \operatorname{pr}_{N,M}\circ
 \mathcal P_{\rad,N}^{\enr}
 =
 \mathcal P_{\rad,M}^{\enr}\circ\pi_{N,M}.
\]
La démonstration n'utilise pas la valeur de la courbe. Elle utilise l'identité des préfixes, qui est l'information que le lecteur réduit à sept coordonnées aurait perdue.

\section{Raffinements non triviaux}

Pour un raffinement \(r:\gamma\to\gamma'\), la naturalité requiert une application \(\widehat r\) sur les registres et l'égalité locale
\[
 h_e
 =
 h_{e'_k}\cdots h_{e'_1}
\]
pour chaque arête \(e\) remplacée par la chaîne raffinée \((e'_1,\ldots,e'_k)\). Si de plus la retenue, le bord et la mémoire de la chaîne raffinée composent les enregistrements de l'étape originale, alors
\[
 \widehat r\circ\mathcal P_{\rad}^{\enr}
 =
 \mathcal P_{\rad}^{\enr}\circ r.
\]
Cette condition est un carré vérifiable ; elle ne doit pas être supposée pour un raffinement arbitraire sans présenter les facteurs.

\section{Équivariance, non invariance}

L'action de \(\Gamma_9\) ramène la phase et prolonge l'état. Dans le codomaine du lecteur agit \(\widehat\Gamma_9^{\rad}\), qui déplace les enregistrements et actualise le dernier préfixe. L'identité correcte est
\[
 \mathcal P_{\rad}^{\enr}(\Gamma_9\gamma)
 =
 \widehat\Gamma_9^{\rad}
 \mathcal P_{\rad}^{\enr}(\gamma).
\]
Exiger l'invariance effacerait précisément la mémoire que la connexion conserve. Dans le secteur réversible, l'inversion des voies satisfait
\[
 H_{C_{\mathrm{rev}}\gamma}=H_\gamma^{-1}
\]
lorsque tous les facteurs appartiennent au sous-groupe affine inversible. Cette affirmation typée ne s'étend pas automatiquement à une involution globale sur n'importe quel quotient ultérieur.

\section{3 démonstrations de non-fidélité}

\begin{proposition}[Cercle intrinsèque et cercle projeté]
Il existe deux histoires ayant la même courbe visible
\[
 c(\theta)=(\cos\theta,\sin\theta)
\]
et des registres distincts.
\end{proposition}
\begin{proof}
La première histoire possède une mémoire axiale constante \(m=0\). La seconde est l'hélice
\[
 h(\theta)=(\cos\theta,\sin\theta,\theta/2\pi),
\]
dont le registre enregistre \(m(\theta)=\theta/2\pi\). Toutes deux satisfont \(\pi_{xy}\circ h=c\). La projection est donc non injective.
\end{proof}

\begin{proposition}[Même extrémité, ordre affine différent]
Si \(a(1-\lambda)\neq0\), les chemins \(A_aM_\lambda\) et \(M_\lambda A_a\) ne sont pas équivalents dans le lecteur enrichi, bien qu'ils puissent être calibrés pour publier une même extrémité.
\end{proposition}
\begin{proof}
Leur commutateur est \(A_{a(1-\lambda)}\neq I\). La donnée du commutateur appartient au registre de préfixes et disparaît lorsqu'on ne conserve que l'extrémité.
\end{proof}

\begin{proposition}[Même courbe, région généalogique différente]
L'égalité de la valeur réelle publiée n'identifie pas les cinq régions de la microfibre de \(\pi\) ; elle ne suffit donc pas non plus à identifier la classe spirale lorsque le lecteur conserve la région, le bord ou la retenue.
\end{proposition}

\chapter{Taxonomie stratifiée des correspondances spirales}

\section{Classification par niveau mathématique}

\begin{longtable}{p{.20\textwidth}p{.30\textwidth}p{.39\textwidth}}
\caption{Niveaux qui ne doivent pas être confondus en une seule notion de courbure.}
\label{espiral:tab:niveles-curvatura}\\
\toprule
Niveau&Invariant&Signification\\
\midrule
\endfirsthead
\toprule
Niveau&Invariant&Signification\\
\midrule
\endhead
APP&\(F(S,P)=1,\ F(P,S)=-1\)&Courbure discrète de l'ordre d'accumulation et de pliage.\\
TRIT&\(\tau\in\{-1,0,1\}\), \(J_\tau^2=-\tau I\)&Régime local elliptique, parabolique ou hyperbolique du transport.\\
TPK&\(\Gamma_9\), retenue, bord, voie et registre&Holonomie de phase et de mémoire sur l'état enrichi.\\
Radial&\(p,a,[\Lambda,C]\)&Caractères qui déterminent la loi radiale avant de reconnaître une courbe.\\
Archimédien&\(r(\theta)\)&Courbe publiée dans \(\R^2\) ; elle peut oublier la généalogie.\\
Frenet&\(\kappa_{\mathrm{F}},\mathcal T_{\mathrm{F}}\)&Courbure et torsion différentielles de la courbe ou de la suspension visible.\\
Dimensionnel&champs, échelles et représentations&Réalisation physique ultérieure, avec ses propres prémisses.\\
\bottomrule
\end{longtable}

\section{Classification par rotation, rayon et mémoire}

\begin{longtable}{ccccl}
\caption{Taxonomie cinématique complète des cas dégénérés minimaux.}
\label{espiral:tab:taxonomia-cinematica-ap}\\
\toprule
\(\Delta\theta\)&\(\Delta u_{\rad}\)&\(\Delta m\)&Salida visible&
Information généalogique\\
\midrule
\endfirsthead
\toprule
\(\Delta\theta\)&\(\Delta u_{\rad}\)&\(\Delta m\)&Salida visible&
Information généalogique\\
\midrule
\endhead
\(0\)&\(0\)&\(0\)&point&état stationnaire publié\\
\(0\)&\(\neq0\)&\(0\)&recta radial&changement d'échelle\\
\(0\)&\(0\)&\(\neq0\)&fibra axial&mémoire sans déplacement transverse\\
\(0\)&\(\neq0\)&\(\neq0\)&feuillet radial suspendu&échelle et mémoire\\
\(\neq0\)&\(0\)&\(0\)&cercle&fermeture angulaire intrinsèque\\
\(\neq0\)&\(0\)&\(\neq0\)&hélice cylindrique&phase fermée, mémoire relevée\\
\(\neq0\)&\(\neq0\)&\(0\)&espiral plana&caractère radial visible\\
\(\neq0\)&\(\neq0\)&\(\neq0\)&suspension hélicoïdale radiale&rotation, rayon et profondeur conservés\\
\bottomrule
\end{longtable}

\section{Classification par holonomie affine}

Le tableau suivant appartient à une réalisation scalaire réelle de la partie linéaire, \(\chi\Lambda=\lambda\chi\) ; pour en alléger la lecture, on écrit \(\Lambda\) dans la colonne des conditions. Il ne classe pas des endomorphismes abstraits au moyen d'inégalités numériques.

\begin{longtable}{p{.18\textwidth}p{.22\textwidth}p{.48\textwidth}}
\caption{Types d'holonomie sur la coordonnée radiale.}
\label{espiral:tab:tipos-holonomia}\\
\toprule
Condition&Orbite de \(u\)&Interprétation\\
\midrule
\endfirsthead
\toprule
Condition&Orbite de \(u\)&Interprétation\\
\midrule
\endhead
\(\Lambda=1,\ C=0\)&\(u_n=u_0\)&rayon constant ; la mémoire peut être nulle ou cachée par projection.\\
\(\Lambda=1,\ C\neq0\)&\(u_n=u_0+nC\)&secteur translationnel qui publie la famille \(L_p\).\\
\(0<\Lambda<1\)&\(u_n\to u_*\)&contraction vers une section fixe.\\
\(\Lambda>1\)&\(|u_n-u_*|\) croît&expansion à partir d'une section fixe.\\
\(\Lambda<0\)&alternance d'orientation&requiert de contrôler le domaine de \(\exp_p\) à chaque étape.\\
\(\Lambda=0\)&\(u_{n+1}=C\)&pliage non inversible ; il appartient au monoïde, non au groupe affine.\\
\bottomrule
\end{longtable}

\section{Critère d'équivalence généalogique}

Deux voies ne sont pas identifiées parce qu'elles publient la même courbe. On exige un isomorphisme typé qui préserve terme à terme :
\[
 \bigl(
 (x_0,H_0,\Led_0);
 (\widehat e_j,p(\widehat e_j),a(\widehat e_j),P_j,A_j,
 h_{e_j},H_j,\Led_j)_{j=1}^{N}
 \bigr)
\]
et qui commute avec les troncatures et raffinements admis. La classification résultante est plus fine que la taxonomie conventionnelle car elle incorpore l'histoire qui précède l'apparition de \(\R^2\).

\chapter{Réalisations ultérieures et contrôles de type}

\section{Mode électromagnétique circulaire}

Soit
\[
 \zeta=\frac{z-c_{\mathrm{vac}}t}{\lambda},
\]
et
\[
 \mathbf E_\sigma
 =E_0(\mathbf e_1\cos2\pi\zeta
 +\sigma\mathbf e_2\sin2\pi\zeta),
 \qquad
 \mathbf B_\sigma
 =c_{\mathrm{vac}}^{-1}\widehat{\mathbf k}\times\mathbf E_\sigma.
\]
Alors
\[
 \mathbf E_\sigma\perp\mathbf B_\sigma\perp\widehat{\mathbf k},
 \quad |E|=c_{\mathrm{vac}}|B|,
 \quad\omega=c_{\mathrm{vac}}k.
\]
À temps fixé, l'extrémité normalisée du champ dans le fibré phase--position décrit
\[
 (\cos2\pi\zeta,\sin2\pi\zeta,\zeta).
\]
Cette affirmation ne dit pas que toute ligne spatiale de champ est une hélice : elle décrit la trajectoire de l'extrémité du vecteur polarisation avec l'événement marqué.

\section{Compatibilité loxodromique}

Pour
\[
 \Phi(r,\varphi,z,t)
 =\ell\varphi+\nu\log(r/r_0)+k_zz-\omega t,
\]
la transformation
\[
 (r,\varphi)\longmapsto
 ((2+\sqrt3)r,\varphi+\pi/2)
\]
préserve la phase modulo \(2\pi\) exactement lorsque
\[
 \nu\log(2+\sqrt3)+\ell\frac{\pi}{2}\in2\pi\Z.
\]
La dilatation s'interprète comme un changement d'échelle entre solutions ; non comme une production spontanée d'énergie.

\section{Centre électronique et lecteur ultérieur}

L'électron occupe l'unique point fixe \((5,5)\) de l'atlas matériel. Cette centralité n'est pas dérivée de nouveau dans le présent travail. La compatibilité opératorielle antérieure comprend
\[
 C^{(3)}m_{\ph}=-\frac12m_{\ph},
 \qquad
 1-\left(\frac12\right)^2=\frac34=\frac{90}{120},
\]
et
\[
 \|[P_{\Sigma,3},P_{\Pi,3}]\|=\frac{\sqrt3}{4}.
\]
La théorie des correspondances spirales ajoute un lecteur de voies ; elle ne remplace pas la Loi générale des masses et n'utilise pas l'électron pour sélectionner rétrospectivement \(p\), \(\Lambda\) ou \(C\).

\section{Matrice fermée de contrôles négatifs}

\begin{longtable}{p{.38\textwidth}p{.50\textwidth}}
\caption{Réfutateurs et erreurs de type exclus par la construction.}
\label{espiral:tab:falsadores-ap}\\
\toprule
Défaillance&Conséquence\\
\midrule
\endfirsthead
\toprule
Défaillance&Conséquence\\
\midrule
\endhead
La courbe classique sélectionne \(p,\Lambda,C\).&Circularité : le lecteur cesse d'être constructif--génératif.\\
Seule l'holonomie finale est conservée.&Perte des préfixes, de la retenue et du bord ; la naturalité n'est pas attestée.\\
L'invariance sous \(\Gamma_9\) est exigée.&L'avancée de mémoire est effacée ; la relation correcte est l'équivariance.\\
On confond \(U_6\), \(w_6\) et \(\F_3^6\).&Confusion d'objets de cardinalités \(468\), \(243\) et \(729\).\\
Le retour de phase est identifié au retour d'état.&Contradiction avec \(B_{K+9}\neq B_K\) en général.\\
On identifie \(\tau\), \(p\) et \(\kappa_{\mathrm{F}}\).&Confusion du régime, du caractère radial et de la courbure visible.\\
Tout rayon variable est qualifié d'hélice générale.&Affirmation différentielle sans vérification de la condition de Lancret.\\
L'axe \(\varphi\) est identifié à l'électron.&Mélange de la fibre des constantes et de la réalisation matérielle.\\
Pell est interprété comme une amplification énergétique.&Violation de la lecture d'échelle et de la conservation physique.\\
L'hélicité électromagnétique est identifiée au spin \(1/2\).&Confusion entre représentations distinctes.\\
Une cellule ponctuelle est imposée à toute espèce.&Contradiction avec le caractère de fibre ou de multisection des espèces non centrales.\\
Les cinq constructions du continu sont séparées.&Rupture de la généalogie unique antérieure à toute publication géométrique.\\
\bottomrule
\end{longtable}

\chapter{Glossaire mathématique}

\begin{longtable}{p{.22\textwidth}p{.66\textwidth}}
\caption{Symboles et termes directeurs.}
\label{espiral:tab:glosario}\\
\toprule
Symbole ou terme&Définition\\
\midrule
\endfirsthead
\toprule
Symbole ou terme&Définition\\
\midrule
\endhead
\(\Sigma,\Pi\)&Évaluations additive et multiplicative sur le support commun APP.\\
\(\tau\)&Régime local d'orientation TRIT ; ce n'est ni le degré radial ni la torsion de Frenet.\\
\(M_{\ph}\)&Sélecteur tritique de phase opératoire.\\
\(\Gamma_9\)&Holonomie nonadique de la dynamique enrichie.\\
\(H_N\)&Holonomie affine accumulée jusqu'au préfixe \(N\).\\
\(\Lambda_N,\mathcal C_N\)&Partie linéaire et cocycle translationnel de \(H_N\).\\
\(p\)&Cocycle entier du caractère radial.\\
\(a\)&Caractère orienté local qui sépare les généalogies de même degré local.\\
\(L_p,\exp_p\)&Coordonnée puissance--logarithmique et son inverse positive.\\
\(\boxplus_p\)&Loi déformée qui publie l'interaction somme--produit.\\
\(\Led_N\)&Registre enrichi de feuillet, retenue, bord, voie et mémoire.\\
\(\mathcal P_{\rad}^{\enr}\)&Lecteur radial enrichi antérieur à la courbe.\\
\(\mathcal P_{\arq}^{\mathrm{esp}}\)&Évaluation ultérieure dans une courbe de \(\R^2\).\\
\(\kappa_{\mathrm{F}},\mathcal T_{\mathrm{F}}\)&Courbure et torsion de Frenet de la courbe visible.\\
Suspension hélicoïdale radiale&Surface ou feuillet à rayon variable qui conserve la phase, la profondeur et le parcours.\\
Publication non fidèle&Évaluation qui peut identifier des histoires distinctes.\\
Constructivo--generativo&Méthode qui construit l'objet et sa généalogie avant de reconnaître son apparence conventionnelle.\\
\bottomrule
\end{longtable}
